\chapter{Anforderungsermittlung}

In der Einleitung wurde deutlich, dass \textit{petappoint} eine fokussierte Lösung für die Terminverwaltung zwischen Tierhaltern und Tierarztpraxen darstellt. Um diese Lösung zielgerichtet umzusetzen, folgt die Anforderungsermittlung dieser Arbeit den Prinzipien des Requirements Engineering. Das ist eine Disziplin, deren Ziel es ist, die Wünsche und Bedürfnisse aller Stakeholder zu verstehen und zu dokumentieren \cite[S.~3]{pohl2021}. Stakeholder sind wiederum Personen oder Organisationen, die Einfluss auf die Anforderungen des Systems haben \cite[S.~2]{pohl2021}. Dazu zählen auch die Nutzenden von petappoint. Eine Anforderung wird dabei als ein notwendiges Bedürfnis eines Stakeholders angesehen, die die Anwendung erfüllen muss \cite[S.~1]{pohl2021}.

Die Analyse der Anforderungen bildet die Grundlage für Konzeption und Entwicklung der vorliegenden Anwendung und Arbeit. Ziel dieses Kapitels ist es, ein vollständiges Bild der Anforderungen zu erarbeiten und diese strukturiert zu dokumentieren.

Zu Beginn wird die Zielgruppe der Anwendung identifiziert und charakterisiert. Darauf aufbauend werden die relevanten funktionalen sowie nicht funktionalen Anforderungen hergeleitet und gegliedert dargestellt. Abschließend werden ausgewählte Anforderungen durch Use Cases konkretisiert und in ihrem praktischen Kontext veranschaulicht.

\section{Zielgruppe}
Die Anwendung richtet sich an zwei primäre Zielgruppen: Tierhalter sowie Tierarztpraxen. Im Rahmen der Analyse bestehender Softwarelösungen in Kapitel 2 wurde deutlich, dass in diesem Bereich keine spezialisierte Terminverwaltungslösung existiert, die beide gleichermaßen adressiert.

Die Anforderungen beider Zielgruppen lassen sich durch folgende Aspekte charakterisieren:

\textbf{Tierhalter}

\begin{itemize}
    \item \textbf{Einfache Bedienbarkeit:} Da die Nutzende in der Regel über keine spezifischen technischen Vorkenntnisse verfügen, muss die Anwendung intuitiv bedienbar sein und ohne umfangreiche Einarbeitung genutzt werden können.
    \item \textbf{Verfügbarkeit rund um die Uhr:} Da berufstätige Tierhalter häufig tagsüber keine Möglichkeit haben, telefonisch einen Termin zu vereinbaren, muss die Buchung zu jeder Tages- und Nachtzeit möglich sein.
    \item \textbf{Haustierzentrierte Verwaltung:} Die Anwendung muss die Anlage mehrerer Haustierprofile mit relevanten Informationen wie Tierart, Rasse und Geburtsdatum unterstützen.
    \item \textbf{Benachrichtigungen:} Automatische Erinnerungen vor bevorstehenden Terminen per E-Mail oder Push-Benachrichtigung sollen vergessene Arztbesuche reduzieren.
    \item \textbf{Datensicherheit:} Da personenbezogene Daten sowie Gesundheitsdaten der Tiere verarbeitet werden, dürfen diese ausschließlich für die jeweiligen Nutzenden zugänglich sein.
\end{itemize}

\textbf{Tierarztpraxen}

\begin{itemize}
    \item \textbf{Digitale Kalender- und Terminverwaltung:} Praxen benötigen ein flexibles System zur Definition von Öffnungszeiten, Terminarten und Kapazitäten, das eine manuelle Terminvergabe per Telefon überflüssig macht.
    \item \textbf{Übersichtliche Tagesplanung:} Ein Dashboard mit Tages- und Wochenansicht aller gebuchten Termine inklusive relevanter Patientendaten ist für einen reibungslosen Praxisablauf erforderlich.
    \item \textbf{Praxisprofil und Auffindbarkeit:} Praxen müssen ihre Leistungen, behandelten Tierarten und Spezialisierungen pflegen können, um von Tierhaltern gezielt gefunden zu werden.
    \item \textbf{Plattformunabhängigkeit:} Die Anwendung soll über gängige Webbrowser nutzbar sein, ohne dass eine lokale Installation erforderlich ist.
\end{itemize}

Zusammenfassend zeigt die Zielgruppenanalyse, dass die Anwendung sowohl für nicht technische Nutzende auf der Seite der Tierhalter als auch für den professionellen Einsatz in tierärztlichen Praxen konzipiert werden muss. Die Berücksichtigung beider Perspektiven ist entscheidend für die Akzeptanz und den nachhaltigen Erfolg des Systems.

\section{Funktionale Anforderungen}

Grundlegend wird zwischen zwei Arten von Anforderungen unterschieden auf die man Einfluss hat: Funktionale Anforderungen und Nicht Funktionale Anforderungen. Funktionale Anforderungen (FA) beschreiben, welche Funktionen und Fähigkeiten das System den Nutzenden zur Verfügung stellen soll \cite[S.~3]{pohl2021}. Zur eindeutigen Formulierung der funktionalen Anforderungen werden 
Satzschablonen eingesetzt. Eine Satzschablone ist ein Bauplan für die 
syntaktische Struktur einer einzelnen Anforderung in natürlicher Sprache 
\cite[S.~71]{pohl2021}. Die funktionalen Anforderungen wurden aus der 
Perspektive beider Zielgruppen abgeleitet:

\textbf{Tierhalter}

\begin{itemize}
    \item \textbf{FA\#01} Die mobile Anwendung \textbf{muss} Nutzenden 
    ermöglichen, sich zu registrieren und anzumelden.

    \item \textbf{FA\#02} Die mobile Anwendung \textbf{muss} Nutzenden 
    ermöglichen, Haustierprofile zu erstellen und zu verwalten.

    \item \textbf{FA\#03} Die mobile Anwendung \textbf{muss} Nutzenden 
    ermöglichen, Tierarztpraxen zu suchen und zu filtern.

    \item \textbf{FA\#04} Die mobile Anwendung \textbf{muss} Nutzenden 
    ermöglichen, verfügbare Termine einzusehen und zu buchen.

    \item \textbf{FA\#05} Die mobile Anwendung \textbf{muss} Nutzenden 
    ermöglichen, gebuchte Termine einzusehen und zu stornieren.

    \item \textbf{FA\#06} Die mobile Anwendung \textbf{muss} Nutzende 
    über Änderungen an ihren Terminen benachrichtigen.
\end{itemize}

\textbf{Tierarztpraxen}

\begin{itemize}
    \item \textbf{FA\#07} Die mobile Anwendung \textbf{muss} 
    Tierarztpraxen ermöglichen, ihr Profil und ihre Verfügbarkeiten 
    zu verwalten.

    \item \textbf{FA\#08} Die mobile Anwendung \textbf{muss} 
    Tierarztpraxen ermöglichen, eingehende Buchungen einzusehen 
    und zu verwalten.
\end{itemize}

\section{Nicht funktionale Anforderungen}

Nicht funktionale Anforderungen sind Qualitätsanforderungen, die sich auf 
die Qualitätsmerkmale des Systems beziehen \cite[S.~4] {pohl2021}.

\begin{itemize}
    \item \textbf{NFA\#01} Die mobile Anwendung muss auf gängigen iOS- und Android-Geräten nutzbar sein.
\end{itemize}

\section{Use Cases}

Use Cases stellen das Verhalten eines Systems aus dem Blickwinkel der Nutzenden dar. Sie beschreiben, wie das System auf eine Interaktion reagiert, welche Ziele damit verfolgt werden und welche alternativen Abläufe dabei auftreten können \cite[S.~15]{cockburn2001}. Aufgrund ihrer textuellen Form eignen sie sich besonders zur klaren und verständlichen Beschreibung funktionaler Anforderungen und ermöglichen eine verständliche Kommunikation ohne besondere fachliche Vorkenntnisse \cite[S.~15]{cockburn2001}. Zur besseren Dokumentation und zur Vollständigkeit werden Sie im Folgenden durch ein Use-Case-Diagramm ergänzt \cite[S.~75]{pohl2021}.

\subsection{Use Case Diagramm}
\begin{landscape}
\begin{figure}[H]
    \centering
    \includesvg[width=0.4\linewidth]{figures/use-case-diagram.svg}
    \caption{Use-Case-Diagramm petappoint}
    \label{fig:usecase-gesamt}
\end{figure}
\end{landscape}

\subsection{Haustier anlegen}

\renewcommand{\arraystretch}{1.5}
\begin{tabular}{|p{5cm}|p{9.2cm}|}
\hline
\textbf{Name} & UC-01: Haustier anlegen \\
\hline
\textbf{Kurzbeschreibung} & Der Tierhalter legt ein neues Haustierprofil 
in der Anwendung an. \\
\hline
\textbf{Akteure} & Tierhalter \\
\hline
\textbf{Vorbedingungen} & 
Der Tierhalter ist eingeloggt. \\
\hline
\textbf{Auslösendes Ereignis} & Der Tierhalter möchte ein neues Haustier 
in der Anwendung erfassen. \\
\hline
\textbf{Hauptszenario} &
1. Der Tierhalter öffnet den Tab „Meine Haustiere". \newline
2. Der Tierhalter tippt auf „Haustier hinzufügen". \newline
3. Die Anwendung öffnet einen neuen Screen mit Eingabefeldern. \newline
4. Der Tierhalter gibt die Daten des Haustieres ein. \newline
5. Der Tierhalter tippt optional auf „Profilbild hinzufügen" und 
wählt ein Bild über die Gerätekamera oder Galerie aus. \newline
6. Der Tierhalter tippt auf „Haustier anlegen". \newline
7. Die Anwendung validiert die Daten. \newline
8. Die Anwendung speichert das Haustier. \newline
9. Der Tierhalter wird zum Tab 'Meine Haustiere' weitergeleitet. \\
\hline
\textbf{Alternativszenario} &
1. Pflichtfelder sind leer $\rightarrow$ Der Button bleibt deaktiviert. \newline
2. Der Tierhalter verweigert den Kamerazugriff $\rightarrow$ Die Anwendung 
zeigt einen Hinweis und erlaubt die Auswahl aus der Galerie. \newline
3. Der Tierhalter verweigert auch den Galerie-Zugriff $\rightarrow$ Das Haustier kann auch ohne Profilbild angelegt werden. \\
\hline
\textbf{Nachbedingungen} &
Das neue Haustier ist unter „Meine Haustiere" sichtbar. \\
\hline
\end{tabular}

\subsection{Haustier bearbeiten}
\renewcommand{\arraystretch}{1.5}
\begin{tabular}{|p{5cm}|p{9.2cm}|}
\hline
\textbf{Name} & UC-02: Haustier bearbeiten \\
\hline
\textbf{Kurzbeschreibung} & Der Tierhalter bearbeitet die Daten eines 
bestehenden Haustierprofils. \\
\hline
\textbf{Akteure} & Tierhalter \\
\hline
\textbf{Vorbedingungen} & 
Der Tierhalter ist eingeloggt. \newline
Der Tierhalter hat mindestens ein Haustier angelegt. \\
\hline
\textbf{Auslösendes Ereignis} & Der Tierhalter möchte die Daten eines 
Haustieres aktualisieren. \\
\hline
\textbf{Hauptszenario} &
1. Der Tierhalter öffnet den Tab „Meine Haustiere". \newline
2. Der Tierhalter tippt auf „Bearbeiten" in der Karte des gewünschten 
Haustieres. \newline
3. Die Anwendung öffnet einen neuen Screen mit den aktuellen Daten 
des Haustieres. \newline
4. Der Tierhalter ändert die gewünschten Daten. \newline
5. Der Tierhalter tippt optional auf „Profilbild ändern" und wählt 
ein neues Bild über die Gerätekamera oder Galerie aus. \newline
6. Der Tierhalter tippt auf „Speichern". \newline
7. Die Anwendung validiert die Daten. \newline
8. Die Anwendung speichert die Änderungen und leitet zurück zu 
„Meine Haustiere". \\
\hline
\textbf{Alternativszenario} &
1. Pflichtfelder sind leer $\rightarrow$ Der Button bleibt deaktiviert. \newline
2. Der Tierhalter verweigert den Kamerazugriff $\rightarrow$ Die Anwendung 
zeigt einen Hinweis und erlaubt die Auswahl aus der Galerie. \newline
3. Der Tierhalter verweigert auch den Galerie-Zugriff $\rightarrow$ Die 
Anwendung zeigt einen Hinweis und das Profilbild bleibt unverändert. \\
\hline
\textbf{Nachbedingungen} &
Das Haustier wird unter „Meine Haustiere" mit den aktualisierten 
Daten angezeigt. \\
\hline
\end{tabular}

\subsection{Haustier löschen}
\renewcommand{\arraystretch}{1.5}
\begin{tabular}{|p{5cm}|p{9.2cm}|}
\hline
\textbf{Name} & UC-03: Haustier löschen \\
\hline
\textbf{Kurzbeschreibung} & Der Tierhalter löscht ein bestehendes 
Haustierprofil aus der Anwendung. \\
\hline
\textbf{Akteure} & Tierhalter \\
\hline
\textbf{Vorbedingungen} & 
Der Tierhalter ist eingeloggt. \newline
Der Tierhalter hat mindestens ein Haustier angelegt. \\
\hline
\textbf{Auslösendes Ereignis} & Der Tierhalter möchte ein Haustierprofil 
dauerhaft entfernen. \\
\hline
\textbf{Hauptszenario} &
1. Der Tierhalter öffnet den Tab „Meine Haustiere". \newline
2. Der Tierhalter tippt auf Bearbeiten im Bereich des gewünschten 
Haustieres. \newline
3. Der Tierhalter tippt auf „Haustier löschen". \newline
4. Die Anwendung öffnet einen Bestätigungs-Bottom-Sheet mit einer 
Warnmeldung. \newline
5. Der Tierhalter bestätigt den Löschvorgang. \newline
6. Die Anwendung löscht das Haustierprofil. \newline
7. Der Tierhalter wird zum Tab „Meine Haustiere" weitergeleitet. \\
\hline
\textbf{Alternativszenario} &
1. Der Tierhalter bricht den Bestätigungs-Bottom-Sheet ab 
$\rightarrow$ Das Haustier bleibt erhalten. \\
\hline
\textbf{Nachbedingungen} &
Das gelöschte Haustier ist nicht mehr unter „Meine Haustiere" sichtbar. \\
\hline
\end{tabular}

\subsection{Termin buchen}
\renewcommand{\arraystretch}{1.5}
\begin{tabular}{|p{5cm}|p{9.2cm}|}
\hline
\textbf{Name} & UC-04: Termin buchen \\
\hline
\textbf{Kurzbeschreibung} & Der Tierhalter bucht einen Termin bei einer 
Tierarztpraxis für sein Haustier. \\
\hline
\textbf{Akteure} & Tierhalter \\
\hline
\textbf{Vorbedingungen} & 
Der Tierhalter ist eingeloggt. \newline
Der Tierhalter hat mindestens ein Haustier angelegt. \\
\hline
\textbf{Auslösendes Ereignis} & Der Tierhalter möchte einen Tierarzttermin 
für sein Haustier vereinbaren. \\
\hline
\textbf{Hauptszenario} &
1. Der Tierhalter tippt auf „Termin buchen". \newline
2. Der Tierhalter gibt optional einen Ort, ein Haustier und eine 
Behandlungsart an. \newline
3. Der Tierhalter tippt auf „Suchen". \newline
4. Die Anwendung zeigt passende Tierarztpraxen in der Nähe in einer 
scrollbaren Liste an. \newline
5. Der Tierhalter tippt auf eine Praxis. \newline
6. Die Anwendung zeigt die Praxisdetails sowie verfügbare Termine 
in einer scrollbaren Kalenderansicht. \newline
7. Der Tierhalter wählt ein Datum und eine Uhrzeit aus. \newline
8. Der Tierhalter tippt auf „Bestätigen". \newline
9. Die Anwendung zeigt eine Buchungsübersicht. \newline
10. Der Tierhalter tippt auf „Termin buchen". \newline
11. Die Anwendung speichert den Termin. \newline
12. Der Tierhalter erhält eine Push-Benachrichtigung als 
Buchungsbestätigung. \newline
13. Der Tierhalter wird zum Tab „Meine Termine" weitergeleitet. \\
\hline
\textbf{Alternativszenario} &
1. Kein Ort angegeben $\rightarrow$ Die Anwendung ermittelt den 
Standort automatisch per GPS und zeigt nahegelegene Praxen. \newline
2. Der Tierhalter verweigert den GPS-Zugriff $\rightarrow$ Die Anwendung 
fordert den Tierhalter zur manuellen Ortseingabe auf. \newline
3. Keine freien Termine verfügbar $\rightarrow$ Die Anwendung zeigt 
einen entsprechenden Hinweis. \\
\hline
\textbf{Nachbedingungen} &
Der gebuchte Termin ist unter „Meine Termine" sichtbar. \\
\hline
\end{tabular}

\subsection{Termin bearbeiten}
\renewcommand{\arraystretch}{1.5}
\begin{tabular}{|p{5cm}|p{9.2cm}|}
\hline
\textbf{Name} & UC-05: Termin bearbeiten \\
\hline
\textbf{Kurzbeschreibung} & Der Tierhalter ändert einen bestehenden 
gebuchten Termin. \\
\hline
\textbf{Akteure} & Tierhalter \\
\hline
\textbf{Vorbedingungen} & 
Der Tierhalter ist eingeloggt. \newline
Der Tierhalter hat mindestens einen gebuchten Termin. \\
\hline
\textbf{Auslösendes Ereignis} & Der Tierhalter möchte Datum, Uhrzeit 
oder Haustier eines Termins ändern. \\
\hline
\textbf{Hauptszenario} &
1. Der Tierhalter öffnet den Tab „Meine Termine". \newline
2. Der Tierhalter tippt auf den gewünschten Termin. \newline
3. Der Tierhalter tippt auf „Termin bearbeiten". \newline
4. Die Anwendung öffnet einen neuen Screen mit den aktuellen 
Termindaten. \newline
5. Der Tierhalter ändert die gewünschten Daten in einer 
scrollbaren Kalenderansicht. \newline
6. Der Tierhalter tippt auf „Speichern". \newline
7. Die Anwendung validiert die Daten. \newline
8. Die Anwendung speichert die Änderungen und leitet zurück 
zu „Meine Termine". \newline
9. Der Tierhalter erhält eine Push-Benachrichtigung als 
Bestätigung der Änderung. \\
\hline
\textbf{Alternativszenario} &
1. Der gewünschte Termin ist nicht mehr verfügbar $\rightarrow$ 
Die Anwendung zeigt eine Fehlermeldung und schlägt alternative 
freie Termine vor. \\
\hline
\textbf{Nachbedingungen} &
Der Termin wird unter „Meine Termine" mit den aktualisierten 
Daten angezeigt. \\
\hline
\end{tabular}

\subsection{Termin stornieren}
\renewcommand{\arraystretch}{1.5}
\begin{tabular}{|p{5cm}|p{9.2cm}|}
\hline
\textbf{Name} & UC-06: Termin stornieren \\
\hline
\textbf{Kurzbeschreibung} & Der Tierhalter storniert einen bestehenden 
gebuchten Termin. \\
\hline
\textbf{Akteure} & Tierhalter \\
\hline
\textbf{Vorbedingungen} & 
Der Tierhalter ist eingeloggt. \newline
Der Tierhalter hat mindestens einen gebuchten Termin. \\
\hline
\textbf{Auslösendes Ereignis} & Der Tierhalter möchte einen gebuchten 
Termin absagen. \\
\hline
\textbf{Hauptszenario} &
1. Der Tierhalter öffnet den Tab „Meine Termine". \newline
2. Der Tierhalter tippt auf den gewünschten Termin. \newline
3. Der Tierhalter tippt auf „Termin stornieren". \newline
4. Die Anwendung öffnet einen Bestätigungs-Bottom-Sheet mit 
einer Warnmeldung. \newline
5. Der Tierhalter bestätigt die Stornierung. \newline
6. Die Anwendung hebt die Buchungsverbindung auf und gibt 
den Termin frei. \newline
7. Der Tierhalter erhält eine Push-Benachrichtigung als 
Bestätigung der Stornierung. \newline
8. Der Tierhalter wird zum Tab „Meine Termine" weitergeleitet. \\
\hline
\textbf{Alternativszenario} &
1. Der Tierhalter bricht den Bestätigungs-Bottom-Sheet ab 
$\rightarrow$ Der Termin bleibt gebucht. \\
\hline
\textbf{Nachbedingungen} &
Der stornierte Termin ist nicht mehr unter „Meine Termine" sichtbar. \\
\hline
\end{tabular}

\subsection{Buchungen einsehen}
\renewcommand{\arraystretch}{1.5}
\begin{tabular}{|p{5cm}|p{9.2cm}|}
\hline
\textbf{Name} & UC-07: Buchungen einsehen \\
\hline
\textbf{Kurzbeschreibung} & Die Tierarztpraxis sieht alle gebuchten 
Termine in einer Tages- und Wochenansicht ein. \\
\hline
\textbf{Akteure} & Tierarztpraxis \\
\hline
\textbf{Vorbedingungen} & 
Die Tierarztpraxis ist eingeloggt. \\
\hline
\textbf{Auslösendes Ereignis} & Die Tierarztpraxis möchte einen 
Überblick über alle anstehenden Buchungen erhalten. \\
\hline
\textbf{Hauptszenario} &
1. Die Tierarztpraxis öffnet den Tab „Kalender". \newline
2. Die Anwendung zeigt eine Tagesansicht mit allen gebuchten 
Terminen des aktuellen Tages. \newline
3. Die Tierarztpraxis tippt auf „Wochenansicht" um zur 
Wochenübersicht zu wechseln. \newline
4. Die Tierarztpraxis tippt auf einen Termin. \newline
5. Die Anwendung zeigt die Detailansicht des Termins mit 
Name des Tierhalters, Tierart, Rasse und Behandlungsart. \\
\hline
\textbf{Alternativszenario} &
1. Keine Buchungen vorhanden $\rightarrow$ Die Anwendung zeigt 
einen entsprechenden Hinweis. \\
\hline
\textbf{Nachbedingungen} &
Die Tierarztpraxis hat einen vollständigen Überblick über 
alle anstehenden Termine. \\
\hline
\end{tabular}

\subsection{Terminslot anlegen}
\renewcommand{\arraystretch}{1.5}
\begin{tabular}{|p{5cm}|p{9.2cm}|}
\hline
\textbf{Name} & UC-08: Terminslot anlegen \\
\hline
\textbf{Kurzbeschreibung} & Die Tierarztpraxis legt einen neuen 
verfügbaren Terminslot in der Anwendung an. \\
\hline
\textbf{Akteure} & Tierarztpraxis \\
\hline
\textbf{Vorbedingungen} & 
Die Tierarztpraxis ist eingeloggt. \\
\hline
\textbf{Auslösendes Ereignis} & Die Tierarztpraxis möchte einen 
neuen Terminslot für Tierhalter verfügbar machen. \\
\hline
\textbf{Hauptszenario} &
1. Die Tierarztpraxis öffnet den Tab „Terminverwaltung". \newline
2. Die Anwendung zeigt eine scrollbare Kalenderansicht mit den 
aktuellen Terminslots. \newline
3. Die Tierarztpraxis tippt auf einen freien Zeitraum. \newline
4. Die Anwendung öffnet einen neuen Screen zur Konfiguration 
des Terminslots. \newline
5. Die Tierarztpraxis gibt Datum, Uhrzeit, Dauer und Behandlungsart an. \newline
6. Die Tierarztpraxis tippt auf „Slot speichern". \newline
7. Die Anwendung validiert die Eingaben. \newline
8. Die Anwendung speichert den Terminslot und aktualisiert 
die Kalenderansicht. \\
\hline
\textbf{Alternativszenario} &
1. Pflichtfelder sind leer $\rightarrow$ Der Button bleibt deaktiviert. \\
\hline
\textbf{Nachbedingungen} &
Der neue Terminslot ist in der Kalenderansicht sichtbar und 
für Tierhalter buchbar. \\
\hline
\end{tabular}

\subsection{Terminslot bearbeiten}
\renewcommand{\arraystretch}{1.5}
\begin{tabular}{|p{5cm}|p{9.2cm}|}
\hline
\textbf{Name} & UC-09: Terminslot bearbeiten \\
\hline
\textbf{Kurzbeschreibung} & Die Tierarztpraxis bearbeitet einen 
bestehenden Terminslot oder sperrt einen Zeitraum. \\
\hline
\textbf{Akteure} & Tierarztpraxis \\
\hline
\textbf{Vorbedingungen} & 
Die Tierarztpraxis ist eingeloggt. \newline
Es existiert mindestens ein Terminslot. \\
\hline
\textbf{Auslösendes Ereignis} & Die Tierarztpraxis möchte einen 
bestehenden Terminslot anpassen oder einen Zeitraum sperren. \\
\hline
\textbf{Hauptszenario} &
1. Die Tierarztpraxis öffnet den Tab „Terminverwaltung". \newline
2. Die Anwendung zeigt eine scrollbare Kalenderansicht mit den 
aktuellen Terminslots. \newline
3. Die Tierarztpraxis tippt auf den gewünschten Terminslot. \newline
4. Die Anwendung öffnet einen neuen Screen mit den aktuellen 
Daten des Terminslots. \newline
5. Die Tierarztpraxis ändert die gewünschten Daten. \newline
6. Die Tierarztpraxis tippt auf „Speichern". \newline
7. Die Anwendung validiert die Eingaben. \newline
8. Die Anwendung speichert die Änderungen und aktualisiert 
die Kalenderansicht. \\
\hline
\textbf{Alternativszenario} &
1. Pflichtfelder sind leer $\rightarrow$ Der Button bleibt deaktiviert. \newline
2. Die Tierarztpraxis möchte einen Zeitraum sperren $\rightarrow$ 
Die Tierarztpraxis tippt auf „Zeitraum sperren", wählt Start und 
Ende aus und bestätigt. Die Anwendung markiert den Zeitraum 
als nicht verfügbar. \\
\hline
\textbf{Nachbedingungen} &
Der Terminslot wird in der Kalenderansicht mit den aktualisierten 
Daten angezeigt. \\
\hline
\end{tabular}

\subsection{gebuchten Termin absagen}
\renewcommand{\arraystretch}{1.5}
\begin{tabular}{|p{5cm}|p{9.2cm}|}
\hline
\textbf{Name} & UC-10: Termin absagen \\
\hline
\textbf{Kurzbeschreibung} & Die Tierarztpraxis sagt einen gebuchten 
Termin ab und der betroffene Tierhalter wird automatisch benachrichtigt. \\
\hline
\textbf{Akteure} & Tierarztpraxis \\
\hline
\textbf{Vorbedingungen} & 
Die Tierarztpraxis ist eingeloggt. \newline
Es existiert mindestens ein gebuchter Termin. \\
\hline
\textbf{Auslösendes Ereignis} & Die Tierarztpraxis möchte einen 
gebuchten Termin absagen. \\
\hline
\textbf{Hauptszenario} &
1. Die Tierarztpraxis öffnet den Tab „Kalender". \newline
2. Die Tierarztpraxis tippt auf den gewünschten Termin. \newline
3. Die Tierarztpraxis tippt auf „Termin absagen". \newline
4. Die Anwendung öffnet einen Bestätigungs-Bottom-Sheet mit 
einer Warnmeldung. \newline
5. Die Tierarztpraxis bestätigt die Absage. \newline
6. Die Anwendung gibt den Terminslot frei. \newline
7. Die Anwendung benachrichtigt den betroffenen Tierhalter 
automatisch per Push-Benachrichtigung und E-Mail. \\
\hline
\textbf{Alternativszenario} &
1. Die Tierarztpraxis bricht den Bestätigungs-Bottom-Sheet ab 
$\rightarrow$ Der Termin bleibt gebucht. \\
\hline
\textbf{Nachbedingungen} &
Der abgesagte Termin ist nicht mehr in der Kalenderansicht sichtbar. \newline
Der betroffene Tierhalter wurde über die Absage informiert. \\
\hline
\end{tabular}